\shyph % slovenské delenie slov pre pdfcsplain
\input opmac
\margins/1 a5 (18,18,16,18)mm % nastavenie formátu A5 a všetky okraje na 18mm
\activettchar"   % nastavuje znak " pre zahájenie a ukončenie verbatim textu
\emergencystretch=2em
%% Potlačenie výskytu sluru; odkomentovať až po odladení výstupu
\overfullrule=0pt
\hyperlinks{}{}  % nastavenie č/b klikacích odkazov
%% https://groups.google.com/g/comp.text.tex/c/SbzyM6ZdVwo/m/HuDeL4vKPgsJ
%% (viď Anil Trivedi, 5. 10. 1993, 13:20:49)
\long\def\Ignore#1\endIgnore{\relax}
\def\endIgnore{\relax} % možné použiť priamo \Ignore ... \endIgnore
\def\IfIgnore{0} % nastaví sa zobrazenie celého (aj doplneného) textu
\def\IfIgnore{1} % po zakomentovaní riadku sa zobrazí doplnený text 

%% nastaví lámanie riadkov medzi \begtt ... \endtt, viď 
%% https://petr.olsak.net/opmac-tricks.html#verbbreak
\def\tthook{\rightskip=0pt plus1fil
  \def\ {\discretionary{\hbox{$\;\setminus$}}{}{\kern.5em}}%
  \everypar={\countspaces}}
\def\countspaces#1\fi{#1\fi\hangindent=\ttindent \countspacesA}
\def\countspacesA{\futurelet\next\countspacesB}
\def\countspacesB{\ifx\next\spacemacro \expandafter\countspacesC\fi}
\def\countspacesC#1{\advance\hangindent by.5em\ \countspacesA}
\def\spacemacro{\ }

%% Nepoužil sa program vlna pre zdrojový text, preto:
\input  encxvlna



%%% Vlastný text %%%

\nonum
\sec Úprava bash skriptu texloop pre beh s xdvi v Ubuntu 20.4 (a novšie)

V pôvodnom skripte texloop z \url{https://petr.olsak.net/texinunix.html}
je pre vetvu xdvi s aktualizáciou prehliadača použitý signál
"kill" "-USR1" "$VIEWPID" (riadok 72, pôvodne riadok 71). V dnešných %$
systémoch UBUNTU, alebo DEBIAN sa však prehliadač nespúšťa priamo, ale cez
wrapper "/usr/bin/perl" "-w" "/usr/bin/xdvi jemny", kde jemny značí meno 
\TeX-ového súboru:

\begtt
$ps -a
4407 pts/2 S 0:00 /bin/bash /home/marian/bin/texloop csplain jemny A
4415 pts/2 S 0:00 /usr/bin/perl -w /usr/bin/xdvi jemny
4419 pts/2 S 0:01 xdvi.bin -name xdvi jemny
\endtt %$

Problém je, že zaslaním signálu USR1 na proces s wrapperom, dôjde nie 
k aktualizovaniu prehliadača xdvi, ale k jeho ukončeniu -- samotný proces 
s prehliadačom ma iné PID. Upravený skript má preto zakomentovaný riadok

\begtt
# kill -USR1 $VIEWPID
\endtt %$
a nad ním je odkaz na podrobný popis problému:

\begtt
# https://bugs.launchpad.net/ubuntu/+source/texlive-bin/+bug/399148
\endtt

V upravenej verzii som použil dirty hack, kde namiesto premennej 
"$VIEWPID" %$
v ktorej sa nachádza PID wrappera sa na riadku 72 nachádza toto:

\begtt
kill -USR1 $(ps a | grep xdvi.bin | awk 'NR==1{print $1}')
\endtt

Výraz v zátvorke má túto funkciu: 
"ps a | grept xdvi.bin |" vypíše všetky procesy, kde sa nachádza reťazec
"xdvi.bin" a pošle ich rúrou programu awk, ktorý vyberie z tohoto zoznamu 
1.~riadok a v prvom stĺpci riadku sa nachádza požadované PID. Problém môže
nastať, ak sa už aj predtým spustil prehliadač "xdvi.bin" -- vtedy však
nenastane ukončenie procesu, len sa náš prehliadač xdvi spustený skriptom
texloop neaktualizuje (na to však stačí myšou ťuknúť do okna s dvi obrazom,
alebo stlačiť "Shift-R" (Reread dvi file), alebo "Ctrl-R" (Redisplays the
current page).



\nonum\secc Natavenia pre xdvi forward a reverse search:

\indent\item
- preklad z tex do dvi robiť s parametrom "-src-specials",%
  \ifnum\IfIgnore = 1 \Ignore \fi
  \fnote{%
  Výstupom je dvi súbor, ktorý sa v niektorých prípadoch {\em môže} líšiť od
  dvi súboru, keby preklad prebehol bez tohto parametra aj v mieste zlomu
  strany, takže tento môže mať iný počet strán. Príkladom je tento text vo
  formáte tex (bez tejto poznámky a následného textu má pri použití veľkosti
  strany A5 dve presne strany, ale ak sa použije "-src-specials" budú to 
  strany tri). To isté platí aj pre pdf súbor. Pre tento formát nemá toto 
  rozšírenie žiadny význam, preto sa jeho použitie v texloop pre formát pdf 
  automaticky potlačí. Pokiaľ chcete použitie "-src-specials" potlačiť aj 
  pre dvi formát, stačí ak posledný parameter skriptu texloop (t. j. tretí
  alebo štvrtý) je "n".}
  \endIgnore
  \ viď riadky 
  23, 24, 62 a 63 texloop; v tomto prípade sa v xdvi po zadaní Ctrl+V
  zobrazia malé červené štvorčeky na začiatku odstavca, kde sa premiestni
  kurzor vim-u v zdrojovom texte po zadaní Ctrl+ľavé tlačidlo myši v xdvi.

\item
- okrem generovania súboru "~/.tlpid" sa generuje aj súbor "~/.tlfname" s 
  názvom zdrojového súboru, aby dochádzalo len k prekladu tohoto súboru aj v
  prípade rozdelenia zdroja do viacerých súborov -- viď riadok 92 texloop

\item
- pre forward search (zo zdrojového tex editovanom vo vim-e do dvi zobrazenom 
  v xdvi) nastavené makro vim-u na klávesu "F3" -- vo ".vimrc" je to táto 
  časť:

\begtt
""" pôvodný variant z https://vim.fandom.com/wiki/Vim_can_interact_with_xdvi
""" (funguje len pre rovnaký nšzov súboru *.tex a *.dvi)
"map <F3> :execute "!xdvi -sourceposition " . line(".") . expand("%") . " " .  expand("%:r") . ".dvi"<CR><CR>
""" názov súboru sa preberá z obsahu súboru ~/.tlfname, ktorý vytvára
""" skript texloop
let fname = readfile(glob('~/.tlfname'))[0]
map <F3> :execute "!xdvi -sourceposition " . line(".") . expand("%") . " " . fname . ".dvi"<CR><CR>
imap <F3> <ESC>:execute "!xdvi -sourceposition " . line(".") . expand("%") . " " . fname . ".dvi"<CR><CR> i
\endtt

\item
- pre reverse search je potrebné spustiť xdvi s definovaním parametra
  "-editor":

\begtt
xdvi -editor "uxterm -e vim --servername xdvi --remote" $MAINFILEBASE &
\endtt %$
viď riadok 46 texloop\medskip

\item
- editor vim je potrebné spustiť s parametrom "vim --servername" "xdvi"; pre
toto mám v "~/.bash_aliases" vytvorený alias pre tvim -- riadok:

\begtt
alias tvim='vim --servername xdvi'
\endtt

Pre takto spustený vim sa bude premiestňovať kurzor z miesta v xdvi do 
zdrojového textu tex otvoreného vo vim. Pokiaľ nie je spustená ani jedna 
inštancia takéhoto vim-u, otvorí sa nové okno UXterm-u a v ňom sa bude presúvať
kurzor na požadované miesto.

\ifnum\IfIgnore = 1 \Ignore \fi
\nonum\secc Výber veľkosti strany pre xdvi

Do pôvodného skriptu bola pri otváraní aplikácie xdvi doplnená aj možnosť
výberu formátu strany -- okrem prednastavenej hodnoty A4 kde možné zadať aj
iné formáty, ktoré pozná aplikácia xdvi (letter, A5, A6, B5, B6, ...)
\endIgnore  
  
\bye
